\documentclass[12pt]{article}
\usepackage[T1]{fontenc}
\usepackage[utf8]{inputenc}
\usepackage{lmodern}
\usepackage{textcomp}
\usepackage{enumitem}
\usepackage{geometry}
\geometry{margin=1in}
\begin{document}
\begin{center}
Answer Key - Chemistry - Graduate Level Assessment 25 \
\small (Answers are ordered exactly as in the question paper.)
\end{center}

\section*{Multiple Choice}
\begin{enumerate}[label=\arabic*.]
\item \textbf{Answer:} C
\\ \textit{Options:} A) 0.3; B) 0.7; C) 0.15; D) 0.9; E) 0.5
\\ \textit{Note:} The mole fraction of H₂SO₄ is calculated by determining the number of moles of H₂SO₄ in the solution and the total number of moles of solute and solvent, where the density and molarity provide the necessary values to find the mass of the solution and thus the amount of water present, leading to a mole fraction of 0.15.
\item \textbf{Answer:} E
\\ \textit{Options:} A) 1.5 \ensuremath{\times} 10^-4 grams; B) 2.50 grams; C) 3.33 grams; D) 4.00 grams; E) 1.00 grams
\\ \textit{Note:} After 2 hours, which is 120 minutes, the material undergoes 3 complete half-lives (36 minutes each), reducing the initial 10.00 grams to 1.25 grams after the first half-life, 0.625 grams after the second, and 0.3125 grams after the third, confirming that the remaining amount after 120 minutes is indeed 1.00 grams.
\item \textbf{Answer:} B
\\ \textit{Options:} A) Redox reactions; B) Elementary processes; C) Dissolution reactions; D) Polymerization reactions; E) Complexation reactions
\\ \textit{Note:} The rate law may be written using stoichiometric coefficients for elementary processes because these reactions occur in a single step, allowing the reaction order to directly correspond to the coefficients of the reactants in the balanced equation.
\item \textbf{Answer:} A
\\ \textit{Options:} A) CH 4; B) NH 3; C) CCl 2 F$_{2}$; D) H$_{2}$O; E) SF 6
\\ \textit{Note:} CH 4, methane, has a single carbon atom bonded to four hydrogen atoms, and does not have any resonance structures due to the absence of pi bonds or lone pairs that would allow for electron delocalization.
\item \textbf{Answer:} A
\\ \textit{Options:} A) 1; B) 0.75; C) 2; D) 4; E) 0.5
\\ \textit{Note:} The uncertainty $\Delta L_z$ for the hydrogen atom in the $2 p_z$ state is 1 because this state has a well-defined angular momentum quantum number $l=1$, resulting in a maximum projection of angular momentum along the z-axis given by the formula $\Delta L_z = m \hbar$ where $m=1$, thus yielding $\Delta L_z = 1 \hbar$.
\end{enumerate}

\section*{True / False}
\begin{enumerate}[label=\arabic*.]
\setcounter{enumi}{5}
\item \textbf{Answer:} True
\\ \textit{Options:} A) True; B) False
\\ \textit{Note:} The statement is true because the average concentration of sodium ions in human blood serum is indeed around 0.15 M, which is critical for maintaining physiological functions such as osmotic balance and nerve impulse transmission.
\item \textbf{Answer:} True
\\ \textit{Options:} A) True; B) False
\\ \textit{Note:} The statement is true because the half-life of ^33 P, which is a radionuclide used in various applications, has been experimentally established and confirmed to be approximately 25.3 days through decay measurements.
\item \textbf{Answer:} False
\\ \textit{Options:} A) True; B) False
\\ \textit{Note:} The statement is false because the decomposition of potassium chlorate (KClO 3) produces oxygen gas (O$_{2}$) in a stoichiometric ratio that would yield more than 0.96 g of oxygen from 4.5 g of KClO 3, specifically around 1.5 g of O$_{2}$ based on the balanced chemical equation.
\item \textbf{Answer:} False
\\ \textit{Options:} A) True; B) False
\\ \textit{Note:} Lead is a solid metal at room temperature, has a silvery-gray color, and is a good conductor of electricity, making the statement false.
\item \textbf{Answer:} False
\\ \textit{Options:} A) True; B) False
\\ \textit{Note:} The pH of a 0.0001 N HCl solution is actually 4, indicating it is acidic, not neutral, as neutral solutions have a pH of 7.
\end{enumerate}

\section*{Long Form Response}
\begin{enumerate}[label=\arabic*.]
\setcounter{enumi}{10}
\item \textbf{Answer:} To determine the relative abundances of the isotopes ^210 X and ^212 X, we can use the equation for the average atomic weight, which is a weighted sum of the isotopes' masses based on their relative abundances. Let the abundance of ^210 X be \( x \) and that of ^212 X be \( 1-x \). The average atomic weight can be expressed as \( 209.64 x + 211.66(1 - x) = 210.197 \). Solving this equation yields \( x \approx 0.7242 \) for ^210 X, or 72.42\%, and \( 1 - x \approx 0.2758 \) for ^212 X, or 27.58\%. Thus, the relative abundances of the isotopes are approximately 72.42\% for ^210 X and 27.58\% for ^212 X, indicating that ^210 X is the more prevalent isotope in this unknown element.
\\ \textit{Note:} correct because it accurately applies the concept of a weighted average to calculate the relative abundances of isotopes based on their masses and the known average atomic weight of the element.
\item \textbf{Answer:} When 1.273 g of silver (Ag) is deposited from the AgNO 3 solution in the first electrolytic cell, we can calculate the number of moles of silver deposited using its molar mass (approximately 107.87 g/mol), yielding about 0.0118 moles of Ag. According to Faraday's laws of electrolysis, the transfer of charge is proportional to the amount of substance deposited, with silver having a valence of +1. Therefore, the equivalent charge required for this process is 0.0118 moles of electrons. In the second cell, copper (Cu) is reduced from CuSO 4, where the reaction also involves the transfer of two moles of electrons per mole of copper deposited. Thus, the amount of copper deposited can be calculated to be 0.273 g, based on the stoichiometric relationship between the deposited silver and copper. This analysis illustrates the principles of charge conservation and the stoichiometric relationships in electrochemical processes.
\\ \textit{Note:} correct because it accurately applies Faraday's laws of electrolysis to calculate the moles of silver deposited, determines the equivalent charge needed for the reaction, and suggests that the deposition of silver in the first cell is coupled with the deposition of copper in the second cell based on stoichiometric relationships between the two reactions.
\end{enumerate}

\vfill
\noindent\textit{End of Answer Key}
\end{document}